\section{The VeriMig Framework}\label{sec:framework}

Fig.~\ref{fig:arch} shows the pipeline. Each ASE object is paired with the checklist of known
differences, translated by a converter, and checked by four gates. We study four converters: a
deterministic rule baseline, an LLM with a plain prompt, an LLM with the checklist added to the prompt, and
the same LLM with a repair loop.

\begin{figure}[t]
\centering
\begin{tikzpicture}[font=\scriptsize, x=1mm, y=1mm,
  box/.style={draw, rounded corners=2pt, align=center, minimum height=7mm, inner sep=2pt},
  gate/.style={box, fill=blue!7, minimum width=13mm},
  arr/.style={-{Stealth[length=1.6mm]}}]
\node[box, minimum width=17mm] (src) at (0,0) {ASE object\\+ checklist};
\node[box, minimum width=19mm] (conv) at (22,0) {Converter\\rules/LLM/repair};
\node[gate] (g1) at (42,0) {G1\\deploy};
\node[gate] (g2) at (59,0) {G2\\residual lint};
\node[gate] (g3) at (59,-14) {G3\\differential};
\node[gate] (g4) at (40,-14) {G4\\reconcile};
\node[box, fill=gray!10, minimum width=22mm] (ref) at (59,-27) {Reviewed reference\\translation};
\node[box, dashed, draw=red!70!black, text=red!70!black, minimum width=50mm] (fb) at (41,12)
  {repair feedback: G1 error text and G2 findings only};
\draw[arr] (src) -- (conv);
\draw[arr] (conv) -- (g1);
\draw[arr] (g1) -- (g2);
\draw[arr] (g2) -- (g3);
\draw[arr] (g3) -- (g4);
\draw[arr] (ref) -- (g3);
\draw[arr, dashed, red!70!black] (g1.north) -- (g1.north |- fb.south);
\draw[arr, dashed, red!70!black] (g2.north) -- (g2.north |- fb.south);
\draw[arr, dashed, red!70!black] (conv.north |- fb.south) -- (conv.north);
\end{tikzpicture}
\caption{VeriMig pipeline. Dashed arrows show the only feedback the repair loop receives; the reference
used by G3 never reaches the converter.}
\label{fig:arch}
\end{figure}

\textbf{G1 (deploy).} The candidate is executed batch by batch in a new SQL Server database after the case
setup. Session options are fixed (\texttt{ANSI\_NULLS}, \texttt{QUOTED\_IDENTIFIER}, and
\texttt{CONCAT\_NULL\_YIELDS\_NULL} on), matching current client defaults.

\textbf{G2 (residual lint).} Eighteen pattern rules detect ASE constructs that should not remain after
translation, such as \texttt{*=}, \texttt{@@sqlstatus}, \texttt{\%1!} placeholders, unbracketed lock hints,
unsigned and \texttt{big*} types, and lock-scheme clauses. Error-level findings fail the gate; deprecated
but valid types are recorded as warnings. Some silent categories, such as the GROUP BY extension and NULL
concatenation, cannot be detected by patterns.

\textbf{G3 (differential).} The reference and the candidate are deployed in separate databases with the
same setup, and every call runs in both. The candidate passes if, for every call, the following match:
all result sets in order (rows are compared as multisets only when a case declares that order does not
matter), with exact values and no rounding; the procedure's return code, captured by wrapping each
single-procedure call; whether an error was raised and its message text; informational messages, where the
case declares them; and the transaction state after the call (\texttt{@@TRANCOUNT} and
\texttt{XACT\_STATE()}). After the last call, the final contents of every user table must also match, so
side effects that no call selects are still checked; \texttt{rowversion} columns are excluded because their
values are database counters. Column names and error numbers are not compared, since a correct translation
may rename an output column or use a different user error number. Every procedure with output parameters
has calls that read them back, which an automated test enforces. Interface metadata, meaning each
procedure's parameter names, types, and output flags and the column types of each result set, is compared
with the reference and reported separately, because a correct translation may legitimately widen a type
(for example, \texttt{unsigned int} to \texttt{bigint}). The comparison shows agreement on the chosen inputs; it is not a proof of
equivalence.

\textbf{G4 (reconcile).} During data movement, VeriMig compares source and target tables. We evaluate four
detectors: row counts; SQL Server's aggregate checksum, \texttt{CHECKSUM\_AGG(BINARY\_CHECKSUM(*))}, which
combines a 32-bit checksum per row into one value by an algorithm Microsoft does not document in detail;
and two per-row SHA-256 hashes. The naive hash joins column values with a separator and writes NULL as a
marker string, which is ambiguous: NULL and the text \texttt{<NULL>} produce the same input, as do
(\texttt{A|B}, \texttt{C}) and (\texttt{A}, \texttt{B|C}). The unambiguous hash writes each column as
\texttt{N;} for NULL or \texttt{V}\emph{length}\texttt{:}\emph{value}\texttt{;} otherwise, converts values
with \texttt{nvarchar(max)} so nothing is truncated, and uses explicit conversion styles for date, time,
and money columns. Distinct rows therefore always produce distinct hash inputs. Rows are compared as a
multiset of (key, hash) pairs, so missing and duplicate keys are also found.

\textbf{Oracle-separated repair.} Algorithm~\ref{alg:agent} gives the repair loop. It receives only the SQL
Server error text from G1 and the remaining ASE constructs from G2, never G3 outcomes. In a real migration
there is no reference for a customer's procedure, and keeping G3 out of the loop prevents the measured
agreement from being tuned against the evaluation signal.

\begin{algorithm}[t]
\caption{Oracle-separated repair loop}\label{alg:agent}
\footnotesize
\begin{algorithmic}[1]
\Require ASE object $s$, model $M$, checklist $T$, budget $K$
\State $h \gets [\textsc{System}(T), \textsc{User}(s)]$
\For{$i = 1$ \textbf{to} $K$}
  \State $c \gets \textsc{Extract}(M(h))$
  \State $(ok_1, e) \gets \textsc{G1}(c)$;\ \ $(ok_2, f) \gets \textsc{G2}(c)$
  \If{$ok_1 \wedge ok_2$} \textbf{break} \EndIf
  \State $h \gets h \Vert [c, \textsc{Feedback}(e, f)]$ \Comment{no G3 signal}
\EndFor
\State \Return $c$ \Comment{then scored once by G1--G3}
\end{algorithmic}
\end{algorithm}
